\subsection{Problem set}

\begin{problem}
    Deformation gradient interpretation
\end{problem}
Given the deformation map $\varphi$:
$$
    \varphi = \begin{bmatrix}
        x \\ y \\ z
    \end{bmatrix} = 
    \begin{bmatrix}
        1.2X \\
        0.8Y \\
        Z
    \end{bmatrix}
$$
Find $\mathbf{F}$ and $\mathbf{J}$, interpret. 
\\\\ \textbf{Answer:} $F$ is defined as:
\begin{pyblock}
from sympy import Matrix, symbols
X, Y, Z = symbols('X Y Z')
ref_coords = Matrix([X, Y, Z])

# Define the deformation map
phi = Matrix(
    [1.2*X,
     0.8*Y,
     1.0*Z])

# Calculate the deformation gradient (F)
F = phi.jacobian(ref_coords)

# Calculate the Jacobian (det F)
J = F.det()
\end{pyblock}
Finally,

$$
    F = \py{ltx_sym_matrix(F)}
$$

$$
    \det{F} = J = \py{J}
$$
Since $\mathbf{J} < 1$, the body is underdoing compression since the stretch in $\mathbf{X}$ is larger than the stretch in the other directions.

\begin{problem}
    Deformation gradient and rotation
\end{problem}
Find $\mathbf{F}$ for a rotation of 30°. Verify that $\det{\mathbf{F}}=1$ and $\mathbf{F}^{T}\mathbf{F}=\mathbf{I}$ 
\\\\ \textbf{Answer:}:
The rotation matrix $\mathbf{R}$ in 2D is defined as:
$$
    \begin{bmatrix}
        \cos \theta && -\sin \theta \\
        \sin \theta && \cos \theta
    \end{bmatrix}
$$

We replace $\theta = 30$°, and compute $\mathbf{F}$.

\begin{pyblock}
import numpy as np


# Define the angle
theta = np.deg2rad(30)

# Define the deformation gradient
F = np.array([
    [np.cos(theta), -np.sin(theta)],
    [np.sin(theta), np.cos(theta)]
])

# Verify that F^t * F = I
verif1 = np.transpose(F) * F

# Verify that det(F) = 1
verif2 = np.linalg.det(F)

\end{pyblock}
We can then verify that:
$$
    \mathbf{F^T F = I = \py{ltx_num_matrix(np.matmul(np.transpose(F),F))}}
$$
and 

$$
    \mathbf{J} = \det{\mathbf{F}} = \py{verif2}
$$
Since rotation is considered rigid body motion, the volume does not undergo deformation which makes since since the jacobian is unity.


\begin{problem}
    Deformation gradient and polar decomposition
\end{problem}
Given 
$$
    \mathbf{F} = \begin{bmatrix}
        1.3 && -0.375 \\
        0.75 && 0.65
    \end{bmatrix}
$$
Compute the polar decomposition (Find $\mathbf{C}$, $\mathbf{U}$,  $\mathbf{V}$ and $\mathbf{R}$).
\\\\ \textbf{Answer:}: From \cref{eq:gradienttransposeproperty} We take the square root of $\mathbf{F^T F}$ to find $\mathbf{U}$. Then we compute the inverse of $\mathbf{U}$ to find $\mathbf{R}$, see \cref{eq:deformationstretch}:
\begin{pyblock} 
import numpy as np
import scipy as sp

# Define the deformation gradient
F = np.array([
    [1.3, -0.375],
    [0.75, 0.65]
])

# Compute C =  F^T F => right Cauchy-Green deformation tensor
C = np.matmul(np.transpose(F), F)

# Compute U = sqrt(C)
U = np.sqrt(C)

# Compute R = F U^-1
R = np.matmul(F, np.linalg.inv(U))

# Compute V = F R^-1
V = np.matmul(F, np.linalg.inv(R))

# Validation using sp.linalg.polar
R2, U2 = sp.linalg.polar(F)

# Find R using iterative approach

Ak = F

Ak1 = 0.5 * (Ak + np.transpose(np.linalg.inv(Ak)))
Ak2 = 0.5 * (Ak1 + np.transpose(np.linalg.inv(Ak1)))
Ak3 = 0.5 * (Ak2 + np.transpose(np.linalg.inv(Ak2)))
    

\end{pyblock}
Giving,
$$
    C = \py{ltx_num_matrix(C)}
$$

$$
    U = \py{ltx_num_matrix(U)} = U2 = \py{ltx_num_matrix(U2)}
$$

$$
    V = \py{ltx_num_matrix(V)}
$$

$$
    R = \py{ltx_num_matrix(R)} = R2 = \py{ltx_num_matrix(R2)} Ak = \py{ltx_num_matrix(Ak)}
$$
Then, the validation of the iterative approach to find $\mathbf{R}$ using 3 steps
$$
    A_{k(0)} = \py{ltx_num_matrix(Ak)}
$$
$$
    A_{k(1)} = \py{ltx_num_matrix(Ak1)}, \epsilon =\py{f"{np.linalg.norm(Ak - Ak1):.4e}"}
$$
$$
    A_{k(2)} = \py{ltx_num_matrix(Ak2)}, \epsilon = \py{f"{np.linalg.norm(Ak1 - Ak2):.4e}"}
$$
$$
    A_{k(3)} = \py{ltx_num_matrix(Ak3)}, \epsilon = \py{f"{np.linalg.norm(Ak2 - Ak3):.4e}"}
$$
The polar decomposition using \textit{sp.linalg.polar} gives the same results as the manual decomposition using \textit{numpy}. Also, the iterative procedure gives quick results without too much effort to find $\mathbf{R}$

\begin{problem}
    Volume change and $J$
\end{problem}
Compute the volume change given
$$
    \mathbf{F} = \begin{bmatrix}
        1.5 && 0 && 0 \\
        0 && 1.2 && 0 \\
        0 && 0 && 1.8
    \end{bmatrix}
$$
\\\\ \textbf{Answer:}: Compute the jacobian ($J$):
\begin{pyblock} 
import numpy as np

# Define the deformation gradient
F = np.array([
    [1.5, 0, 0],
    [0, 1.2, 0],
    [0, 0, 1.8]
])

# Compute J=det(F)
J = np.linalg.det(F)
\end{pyblock}
Since $J = \py{J} > 1$. The body becomes $\py{J}$ larger, thus the volume is $\py{J}$ times bigger.